% =====================================================================
%  前置部分：译校说明 + 作者前言（まえがき，原书 p.i–ii）
% =====================================================================
\chapter*{译校说明}
\addcontentsline{toc}{chapter}{译校说明}
\markboth{译校说明}{译校说明}

\section*{一、底本}

本书译自植田一石《数物系のためのシンプレクティック幾何学入門》，SGC ライブラリ 118，サイエンス社，2015 年 4 月。原书正文 175 页，另附前言、后记、参考文献与索引。原书以物理学与数学专业的本科生、研究生为主要读者，行文保持讲义式的平易口吻。

中译本以该版 PDF（217 页，含封面、目次等前置页）为底本，逐页对照扫描原图翻译；原文的公式与陈述编号、图表编号一律照录。

\section*{二、体例}

\begin{enumerate}[label=\arabic*.,leftmargin=2.2em,itemsep=2pt]
  \item \textbf{页码对照}：正文中以 \texttt{\% [原文 p.N]} 标注各段所对应的原书页码；图注末尾标「（原书 p.N）」。原书页码 = PDF 页码 − 9。
  \item \textbf{公式与陈述编号}：原书公式按「章.节.序号」编号（如 (2.1.19)、(8.4.33)），各类陈述（定義、定理、命題、補題、系、例、注意、問題、予想、反例、記法、仮定 等）亦按「章.节.序号」编号（如「定義 2.1.1」「命題 3.1.4」）。中译本一律照录原编号，不改动、不重排。
  \item \textbf{图}：原书插图均自扫描页裁出，图号与图注照原书；原书仅 5 幅编号插图（図 2.1–2.4、図 9.1）。
  \item \textbf{译校批注}：以源码注释 \texttt{\% [Tutor Note: ...]} 形式给出，不进入正文排版，全部汇集于书末「附录 A：译校批注」，按章排列并注明原书页码。批注只写三类内容：背景（人物、史事、术语来源、与物理的对应）、难点（推导跳步处的点拨）、校勘（原书疑似笔误及依据）。原书疑似笔误一律照原样排，只在批注中指出。
  \item \textbf{本节重点与章末回顾}：每节末尾的「本节重点」与每章末尾的「本章回顾」方框均为中译本所加，非原书内容。全书可用 \texttt{main.tex} 中的 \texttt{\textbackslash showtakeawayfalse} 与 \texttt{\textbackslash showchapterreviewfalse} 一键隐藏。
  \item \textbf{索引}：原书卷末索引按日文五十音排列并给出日文版页码，于中译本无对应页码可循，故不照录；改以书末「附录 B：术语对照表」代替，列出原书术语、中译名与英文原名。
  \item \textbf{参考文献}：照录原书条目，不译书名，仅对条目编号与作者名的排印错误作最小订正。
\end{enumerate}

\section*{三、术语原则}

术语以「附录 B：术语对照表」为准。原则是：中文学界有稳定、无歧义通行译名者用中文；无稳定译名或硬译反碍对照文献者保留原文（如 operad、Fukaya 范畴、Floer 同调、Gromov–Witten 不变量等专名与缩写）。人名一律不译。日文与中文用法有别者按中文习惯处理，例如：日文「解析力学」译「分析力学」，「多様体」译「流形」，「概複素構造」译「近复结构」，「正則」在复几何语境译「全纯」，「固有値」译「特征值」，「コンパクト」译「紧」，「準同型」译「同态」，「余随伴軌道」译「余伴随轨道」。

\section*{四、声明}

本书为学术研读用译本，不作商业用途。译文与批注中的错误由译者负责，恳请读者指正；原书的一切著作权归作者与サイエンス社所有。

\clearpage

% =====================================================================
\chapter*{前言}
\addcontentsline{toc}{chapter}{前言（原书 p.i–ii）}
\markboth{前言}{前言}

本书主要面向学习物理学与数学的本科生和研究生。不过也设想有意愿强烈的高中生与其他领域的研究者阅读，因此尽量以平易的方式介绍辛几何的标准话题。

说明事物时的顺序大致可分为三种：历史的顺序、逻辑的顺序，以及教育的顺序。历史的顺序既不是抵达目的地的最短路径，又常伴有并不显然的跳跃，要理解其必然性未必容易。逻辑的顺序条理整饬、没有冗余，但定义与定理背后层层积累的智慧被高度抽象化所掩盖，初学者乍看之下会在枯燥的符号罗列中迷惑徘徊。另一方面，教育的顺序则无视历史顺序或逻辑顺序，按心理负担较小、容易信服的顺序展开论述。

本书尝试既不用历史顺序、也不用逻辑顺序，而以教育的顺序展开论述。不过，什么样的议论作为动机说明具有说服力，因人而异、差别颇大，因此本书的议论不过是教育顺序的一个例子。建议读者按各自的兴趣与理解，自行重构本书的内容。

本书的一部分基于笔者在国内外大学与研究所所作的讲义与演讲，但并未改写成教科书体，而是努力保持讲课那种平易的气氛。本书内容从前向后缓慢地相互关联，但为使各章大致可以独立阅读，对重要概念不厌其烦地反复说明。又为使自学的读者不致因简单的事卡住，计算过程写得近乎啰嗦。越往后所需的预备知识越多，但即使遇到不懂之处，也不要就此放弃，可参阅其他文献补足，或者视情况把那里搁置，暂且先往下读。

本书的大致结构如下。第 1 章叙述辛几何的历史与现状，以及各章的背景。这部分是故事，不需要预备知识。第 2 章沿教育的顺序，从分析力学的初等部分走到辛几何，并尽量把所需的预备知识压到最少。若只举出这一章的一个目标，那大概就是理解关于对称性与守恒律的 Noether 定理的主张与证明。Noether 定理是值得细细品味的深刻定理，仅为了解它就值得学习分析力学。对称性可以改写为群作用，构成本书（或者辛几何本身）的主题之一。

若再举出第 2 章的另一个目标，那便是以能让人信服其自然性的方式给出辛流形的定义。不过，数学中的定义乃是经过层层历史的打磨、作为其结果而成为现在的形式的（拓扑空间的定义即是典型），归根结底只能由好的具体例来正当化。

如果说第 2 章是从物理学视角引入辛几何，那么第 3、4 章便是从数学立场引入辛几何。第 3 章从 G 结构几何学的视角重新给出辛流形的定义，第 4 章举出若干辛流形的例子。到此为止是绪论，自第 5 章起进入本论。

第 5 章介绍辛几何中处理群作用的框架。在这里，读者将看到向量场、Lie 群等流形论中所学的各种概念生动活跃的样子。

辛几何在某种意义上作为分析力学的终极表述而诞生，但力学问题并非列出运动方程便告结束；即便求得运动方程，若对其解一无所知，也几乎没有意义。关于解能知道多少，因系统而异、差别巨大，而把可能想到的最好情形表述出来的，便是完全可积系统的概念。完全可积系统的一般理论在 Liouville–Arnold 定理中大体已尽，但其具体例构成代数、几何、分析交汇的丰富世界，研究者络绎不绝。第 6 章就完全可积系统，介绍 Liouville–Arnold 定理的概要，以及典型具体例 Toda 格子，并采取 Lie 群分解的立场。

可积系统是极为例外的存在；对大多数运动方程，都无法用熟知的函数显式表示其解。即便在这种情形下，为了对解得到某种有意义的结果，Poincaré 创立了动力系统的理论。第 7 章简要介绍处理接近可积系统的动力学系统的 KAM 理论，以及处理远离可积系统的动力学系统的遍历理论。前者以太阳系的稳定性、后者以统计力学为动机而产生，两者都是与物理学渊源深厚的理论。

第 8 章处理称为辛环面流形的一类非常特殊的辛流形。粗略地说，它被定义为具有完全可积系统结构的紧辛流形中最好（或最单纯）者。这一领域的基本定理是 Delzant 定理：辛环面流形与满足适当条件的凸多面体一一对应。借助这一定理，关于辛环面流形的问题在原理上全部归结为凸多面体的组合论。

第 9 章处理二维规范理论。二维规范理论与辛几何的密切关系由 Atiyah–Bott 发现，并成为其后发展的指导原理。本章的目标是介绍 Narasimhan–Seshadri 定理的 Donaldson 证明：Riemann 面上的全纯向量丛满足稳定，与它容许直和不可分解的 Yang–Mills 联络，二者等价。这一定理的高维化，以及从规范理论向重力的推广，曾是微分几何的中心问题之一。

第 10 章转换方向，介绍复几何中的辛结构。这与四元数关系密切，形成一个与实辛几何面貌相当不同的独特世界。紧的例子极少，但不紧的例子常常作为场论与超弦理论中的模空间出现。

按最初的计划，除上述内容之外本还打算讨论若干仍在活跃研究中的前沿话题，但因篇幅所限遗憾未能实现。本书未能介绍的此后发展，将在后记中稍作涉及。

承蒙编辑、同事、恩师、家人等诸多人士的帮助，这本书才得以问世，在此表示感谢。

\vspace{1em}
\noindent\hfill 2015 年 4 月\qquad 植田一石\hspace{2em}
